Jean Sibelius (1865–1957) was a Finnish composer whose music played a vital role in shaping national identity during Finland’s struggle for independence from Russia. Best known for his symphonies, tone poems, and the stirring Finlandia, Sibelius developed a distinctive musical voice marked by rich orchestration, organic development, and deep ties to nature and Finnish folklore. His seven symphonies trace a bold evolution in form and harmony, influencing 20th-century composers across Europe. Though he withdrew from composition later in life, his legacy endures as a cornerstone of Nordic music and a symbol of Finnish cultural pride and artistic independence. 

\vspace{\titleheadsep}
\textit{JEL}: D40, D82, L11.\\
\textit{Keywords}: Keyword 1, Keyword 2, Keyword 3.
